% Lösungen Einheit 4 von 4 – Zuordnungstypen erkennen und anwenden (zusammenbau v0.1)
\einheitenkopf{Einheit 4 von 4 · Zuordnungstypen erkennen und anwenden}

\erg{23}{a) nein – die Grundgebühr kostet schon 4 €.}
%% TODO Lösung der Erklärzeile 24g) fehlt in der Bank
\erg{24}{a) mehr \quad b) proportional – alle Quotienten sind 3. \quad c) nicht proportional – 5 : 1 = 5, aber 9 : 2 = 4,5. \quad d) proportional – jeder y-Wert ist zwei Drittel des x-Werts. \quad e) nicht proportional – 7 : 2 = 3,5, aber 17 : 5 = 3,4. \quad f) proportional – alle Quotienten sind 2,5. \quad g) \ldots \quad h) antiproportional – alle Produkte sind 18. \quad i) nicht proportional – schon 0 Stunden kosten 3 €. \quad j) proportional – Gerade durch den Ursprung. \quad k) antiproportional – alle Produkte sind 40. \quad l) Nein: 8 : 2 = 4, aber 14 : 4 = 3,5; die Quotienten sind verschieden.}
%% TODO Lösung der Erklärzeile 25g) fehlt in der Bank
\erg{25}{a) 1,79 € je Liter \quad b) 4 · 2,50 € = 10 € \quad c) 7 · 3 € = 21 € \quad d) 6 · 15 € = 90 € \quad e) 3 · 1,50 € = 4,50 € \quad f) 500 · 0,90 € = 450 € \quad g) \ldots \quad h) 6 : 1,50 = 4 kg \quad i) 45 : 3 = 15 km/h \quad j) 150 : 50 = 3 h \quad k) 30\,000 : 100 · 8 = 2\,400 l; 2\,400 · 1,625 € = 3\,900 € \quad l) 3 km : 4 km/h = 0,75 h = 45 min; 9:50 Uhr + 30 min + 45 min = 11:05 Uhr}
\erg{26}{a) Alter → Körpergröße: Ältere Kinder sind meist größer, aber doppeltes Alter heißt nicht doppelte Größe.}
\erg{27}{a) Lukas hat antiproportional gerechnet; bei gleicher Geschwindigkeit ist Strecke → Zeit proportional: 180 : 90 = 2 h. \quad b) Je mehr Personen teilen, desto weniger zahlt jede – die Werte fallen, bei proportionalen Zuordnungen wachsen beide. \quad c) proportional; Ursprungsgerade durch (0|0) und (4|20) \quad d) 280 : 80 = 3,5 h; Abfahrt spätestens um 10:30 Uhr.}
